\section{Описание применения генеративной модели}
\label{sec:genai}

\textbf{Название модели.} Claude (разработчик~--- Anthropic),
семейство больших языковых моделей.

\textbf{Адрес в сети «Интернет».} \url{https://claude.ai}
(разработчик~--- \url{https://www.anthropic.com}).

\textbf{Цели применения.} Генеративная модель применялась как
вспомогательный инструмент (ассистент программиста и редактор), но не как
источник идей или результатов работы. Основные цели~--- ускорение
рутинных операций написания и отладки кода, техническая редактура и
оформление текста отчёта в \LaTeX{}, а также дополнительная проверка
корректности экспериментов.

\textbf{Способ применения.} Применение носило характер диалогового
ассистирования по конкретным задачам и охватывало:

\begin{itemize}
    \item написание и отладку отдельных фрагментов программного кода
          (парсеры источников данных, реализации моделей HAR, XGBoost и
          kNN, модуль метрик, веб-интерфейс) с последующей проверкой,
          тестированием и адаптацией кода автором;
    \item техническую редактуру и оформление текста отчёта в \LaTeX{};
    \item дополнительную проверку экспериментальной части, в ходе которой
          были подтверждены и устранены методологические неточности
          ранних версий кода (несогласованность формул QLIKE между
          моделями, утечка данных из-за одновременности sentiment-признаков
          с целевой переменной).
\end{itemize}

\textbf{Мой собственный вклад.} Постановку задачи и формулировку цели,
выбор предметной области и источников данных, проектирование архитектуры
программного комплекса, выбор и обоснование моделей и метрик, организацию
схемы оценивания, фактический сбор данных, запуск всех вычислительных
экспериментов, интерпретацию результатов и все выводы работы я выполнил
самостоятельно. Все приведённые в работе численные показатели получены
в результате моих реальных вычислительных экспериментов на собранных
данных. Сгенерированные моделью фрагменты кода и текста я рассматривал
как черновой материал: проверял, тестировал, корректировал и принимал их
под свою ответственность, неся полную ответственность за итоговое
содержание работы.
